%------------------------------------------------------------------------------%
\begin{exercises}

% Ex 28 pg 10
%--------------------------------------------------------------------------------------------------%
\begin{exercise}
  A sequência $(a_n)$ converge ou diverge? Encontre seu limite se for convergente.
  \[
  a_n = \frac{n+(-1)^n}{n}
  \]
  \begin{answer}
    Tentando calcular o limite temos
    \begin{align*}
      \lim_{n\to\infty} a_n & = \lim\frac{n+(-1)^n}{n}                        \\
      & = \lim\left[\frac{n}{n}+\frac{(-1)^n}{n}\right] \\
      & = \lim\left[1+\frac{(-1)^n}{n}\right]
    \end{align*}
    A sequência constante igual a 1 converge, temos que verificar se a sequência $\sfrac{(-1)^n}{n}$
    também converge. Observe que para todo $n$ temos
    \[
    -\frac{1}{n} \;\leq\; \frac{(-1)^n}{n} \;\leq\; \frac{1}{n}
    \]
    Como $\sfrac{-1}{n}$ e $\sfrac{1}{n}$ convergem para zero, pelo teorema do confronto
    para sequências temos que $\sfrac{(-1)^n}{n}$ também converge para zero.
    Podemos, então, escrever
    \begin{align*}
      \lim_{n\to\infty} a_n & = \lim\left[1+\frac{(-1)^n}{n}\right] \\
      & = \lim 1+ \lim \frac{(-1)^n}{n}       \\
      & = 1 + 0 = 1
    \end{align*}
  \end{answer}
\end{exercise}

% 23-56 - Stewart pt pg 633
%------------------------------------------------------------------------------%
\begin{exercise}
  Verifique se a sequência converge ou não e calcule o limite das sequências convergentes.
  \begin{mathlist}[2]
    \item a_n = \frac{(-1)^{n-1}\;n}{n^2+1}         % 33
    \item a_n = \frac{(-1)^n\;n^3}{n^3+2n^2+1}      % 34
    % 41 - 50
    \item a_n = \frac{\cos^2(n)}{2^n}
    \item a_n = 2^{-n} \cos(n\pi)
    \item a_n = \frac{\sin(2n)}{1+\sqrt{n}}
  \end{mathlist}
\end{exercise}

%------------------------------------------------------------------------------%
%\begin{exercise}
%\begin{answer}
%\end{answer}
%\end{exercise}

%------------------------------------------------------------------------------%
%\begin{exercise}
%\begin{mathlist}[2]
%  \item
%  \item
%  \item
%\end{mathlist}
%\begin{answer}
%  \begin{alphalist}[3]
%    \item
%    \item
%    \item
%  \end{alphalist}
%\end{answer}
%\end{exercise}

\end{exercises}
%------------------------------------------------------------------------------%
